Este proyecto se desarrolla en el contexto de la asignatura \textbf{Electrónica de Consumo (ELCO)} de la \textbf{Universidad Politécnica de Madrid (UPM)}, bajo un enfoque de diseño propio de una pequeña empresa de ingeniería.

El objetivo ha sido diseñar y construir una máquina capaz de \textbf{barajar y repartir cartas de forma automática}, integrando sistemas electrónicos, mecánicos y de programación para eliminar la intervención humana y mejorar la rapidez, fiabilidad y repetibilidad del proceso.

Para ello, se ha utilizado una \textbf{Arduino MEGA 2560 R3} como unidad de control, junto con un sistema de motores gestionados mediante un driver \textbf{L298N}. El diseño mecánico se ha realizado mediante piezas impresas en 3D, incluyendo rodillos de arrastre optimizados para el manejo de cartas.

El sistema incorpora además sensores como una \textbf{barrera infrarroja} para la detección de cartas y un \textbf{sensor HOME} para la calibración de posición inicial, así como un \textbf{Slip Ring} que permite la rotación continua del mecanismo sin problemas de cableado.

En conjunto, el proyecto permite aplicar de forma integrada conocimientos de electrónica, mecánica y programación en un sistema funcional de automatización.